\section{Measurable causal graphs: proofs}
\label{sec:annex-p16-measurable-causal-graphs}

\begin{annexitem}{annex:p16.measurable-ci-calculus}{thm:p16-measurable-ci-calculus}{Conditional expectation and one common measure identity}
\begin{proof}
Suppose first that the product-test definition holds. For bounded
\(u(X)\), put \(h(Z)=E[u(X)\mid Z]\). On a generating event
\(\{Y\in D,Z\in E\}\), conditional pullout and the definition give
\[
\begin{aligned}
 E[\mathbf1_E(Z)\mathbf1_D(Y)h(Z)]
 &=E[\mathbf1_E(Z)h(Z)E[\mathbf1_D(Y)\mid Z]]\\
 &=E[\mathbf1_E(Z)u(X)\mathbf1_D(Y)].
\end{aligned}
\]
A pi-lambda argument extends this integral identity to all events in
\(\sigma(Y,Z)\), proving that \(h(Z)\) is a conditional expectation
given \((Y,Z)\). Conversely if the conditional expectation given
\((Y,Z)\) is \(h(Z)\), tower and pullout give
\[
 E[u(X)v(Y)\mid Z]
 =E[h(Z)v(Y)\mid Z]=h(Z)E[v(Y)\mid Z].
\]
The requisite tower and pullout results are
\cref{thm:p15-conditional-tower,thm:p15-conditional-pullout}.

Choose regular kernels for \(X,Y,(X,Y)\) given \(Z\). Their existence
and uniqueness as measures follow from
\cref{thm:p15-value-conditional-law}. Each standard Borel target has
a countable generating pi-system, by its Borel real coding. For every
rectangle from the two countable systems, conditional independence of
its indicators equates the joint kernel to the product of its marginal
kernels outside a null set of \(Z\). Remove the countable union of these
sets. At every remaining \(z\), measure uniqueness gives equality on
the entire product sigma-algebra. The product kernel is measurable:
its values on rectangles are products of measurable marginal values,
and a pi-lambda argument extends this to all measurable sets. Conversely
the common measure identity integrated against products of bounded
functions gives conditional independence. This proves the second
equivalence without separate null-set choices for uncountably many events.

Symmetry is immediate from the defining product identity. Decomposition
follows by taking a test function that ignores an extra coordinate.
For weak union, the no-extra-information identity gives
\(E[u(X)\mid Y,W,Z]=E[u(X)\mid Z]\). The right side is measurable
from \((W,Z)\); tower makes it also \(E[u(X)\mid W,Z]\).
Applying the equivalence in reverse gives \(X\perp Y\mid(W,Z)\).
For contraction the two premises give successively
\[
 E[u(X)\mid Y,W,Z]=E[u(X)\mid Y,Z]=E[u(X)\mid Z],
\]
which is precisely the conclusion. All tests are bounded, so every
conditional expectation used is integrable. Measurable transformations
of independent tuples follow either by composing bounded tests or,
when the functions also depend on \(Z\), by pushing forward the
conditional product kernels at each \(z\). No intersection axiom or
strictly positive law entered any of these steps.
\end{proof}
\end{annexitem}

\begin{annexitem}{annex:p16.measurable-kernel-construction}{thm:p16-measurable-kernel-construction}{Associative kernel combination and legal adjacent swaps}
\begin{proof}
For the empty graph its point-mass law gives all assertions directly.
At the first vertex of a nonempty graph the parent tuple is empty, so its kernel is a
probability measure. At each later stage extend its kernel to the whole
preceding tuple by the measurable parent projection. The kernel
joint-measure theorem constructs a unique new probability measure and
identifies its conditional law of the new coordinate. Induction gives
existence, normalization, the iterated-integral identity, and the
conditional-prefix statement. All parameter integrals are measurable by
the same theorem. Nonnegative tests pass through each step by monotone
convergence; signed tests are handled only after their absolute integrals
are finite. In particular no multiplication of hypothetical densities
has occurred.

For order independence, take two consecutive vertices \(v,w\) in a
topological order with neither a parent nor an ancestor of the other.
Both parent tuples are functions of the common earlier tuple \(z\).
Conditional on \(z\), their two update kernels combine as the product
\(K_v(z,dx_v)\otimes K_w(z,dx_w)\). On rectangles this follows
directly from the kernel joint-measure definition; measure uniqueness
extends it to every measurable set. Fubini for these two probability
measures permits either integration order. Integrating all later
updates first gives a measurable nonnegative function of \((z,x_v,x_w)\),
so the same swap preserves the full law, not just this pair's marginal.

Any two topological orders are connected by such adjacent swaps. To see
this, take the first vertex in the desired order and move it leftward
through the current order. A crossed vertex cannot be its ancestor,
since it would then have to precede it in the desired order; it cannot
be its descendant, since it precedes it in the current order. Thus each
swap is of incomparable vertices and is legal. Fix that first position
and repeat on the remaining finite graph. Induction proves the claim
and hence order independence. This is the source's kernel-combination
argument with the needed product/Fubini domains made explicit.

If \(D\) is ancestral, all parents of its vertices lie in \(D\).
List a topological order of \(G[D]\) first, followed by one of the
remaining vertices. No edge can enter \(D\) from its complement,
so this is a topological order of \(G\). Its initial joint law is
the kernel construction on \(D\), and all later extensions preserve
that marginal. Order independence now proves the asserted ancestral
marginal for every original ordering. Topological prefixes are
ancestral. Their conditional-coordinate identities therefore hold
also in the full law, since conditional expectations from preceding
coordinates depend only on the corresponding joint marginal.
\end{proof}
\end{annexitem}

\begin{annexitem}{annex:p16.measurable-kernel-randomization}{thm:p16-measurable-kernel-randomization}{A jointly measurable quantile and Borel image repair}
\begin{proof}
Choose a measurable isomorphism \(\phi:S\to E\), where \(E\) is a
nonempty Borel subset of \(\mathbb R\), using the proved standard
Borel coding. Push each \(K(t,\cdot)\) through \(\phi\) and regard
it as a probability measure \(\widetilde K_t\) on \(\mathbb R\)
concentrated on \(E\). Its CDF is
\(F_t(x)=\widetilde K_t(( -\infty,x])\). For fixed \(x\), this
is measurable in \(t\). Define the source's upper quantile by the
equivalent rational formula
\[
 q(t,u)=\inf\{r\in\mathbb Q:F_t(r)>u\},\qquad 0<u<1.
\]
For each \((t,u)\), the set is nonempty by the upper CDF limit and
is bounded below by the lower CDF limit, so the infimum is real.
Its strict sublevel sets satisfy
\[
 \{(t,u):q(t,u)<x\}
   =\bigcup_{r\in\mathbb Q,\ r<x}\{(t,u):F_t(r)>u\}.
\]
This countable union is measurable, proving joint measurability.

For fixed \(t,x\), monotonicity and right continuity of the CDF give
\[
 \{u:u<F_t(x)\}\subseteq\{u:q(t,u)\leq x\}
                  \subseteq\{u:u\leq F_t(x)\}.
\]
For the first inclusion use rational arguments arbitrarily close above
\(x\). For the second, if the infimum is at most \(x\), every rational
argument strictly above \(x\) has CDF at least \(u\), and then let
such arguments decrease to \(x\). The two bounding sets have the same
uniform measure \(F_t(x)\). Thus \(q(t,U)\) has CDF \(F_t\), and
has law \(\widetilde K_t\). In particular \(q(t,U)\in E\) almost
surely for each fixed \(t\).

Choose \(s_0\in S\) and define
\[
 F(t,u)=\begin{cases}\phi^{-1}(q(t,u)),&q(t,u)\in E,\\
                         s_0,&q(t,u)\notin E.\end{cases}
\]
The exceptional set is jointly measurable because \(E\) is Borel.
The inverse coding and this measurable default make \(F\) jointly
measurable on its whole domain. For each fixed parameter the modification
is on a uniform-null set, and so its pushforward law is exactly
\(K(t,\cdot)\). This proves the kernel realization at every parameter,
including parameters absent from an observed law.

Apply this construction separately to each of the finitely many DAG
kernels and use independent uniform variables on their finite product
space. Topological evaluation and successive conditioning show that its
joint law is the kernel law. The representation is not unique: a fair
response independent of binary treatment can be represented as
\(Y=U\), or as \(Y=A\oplus U\), with a fair independent bit \(U\).
Both give the same controlled response kernel at both actions, but
their potential responses are respectively equal or opposite. Thus
the construction provides a structural realization, not an inference
of a unique cross-world coupling from a probability kernel.
\end{proof}
\end{annexitem}

\begin{annexitem}{annex:p16.measurable-global-markov}{thm:p16-measurable-global-markov}{The complete density-free terminal-vertex induction}
\begin{proof}
\paragraph{Recursive kernels imply the global property.}
We first prove a full-partition assertion by induction on the number
of vertices. Suppose \(N=D\mathbin{\dot\cup}E\mathbin{\dot\cup}C\)
and there is no edge between \(D\) and \(E\) in the moral graph
after deleting \(C\). We claim \(X_D\perp X_E\mid X_C\).
It is immediate if either side is empty, including the empty and
one-vertex starting cases.

Take the last vertex \(v\) in a topological order. It is terminal.
The remaining law is its kernel-model marginal, and the moral graph
on the remaining DAG has no more edges than the old moral graph
with \(v\) deleted. If \(v\in D\), its parents lie in
\((D\setminus\{v\})\cup C\), since a parent in \(E\) would be
a forbidden moral edge. The conditional kernel of \(X_v\) given
all other variables consequently depends only on
\(X_{D\setminus\{v\}},X_C\). The no-extra-information
characterization gives
\(X_v\perp X_E\mid X_{D\setminus\{v\}},X_C\).
Induction gives \(X_{D\setminus\{v\}}\perp X_E\mid X_C\).
Symmetry and contraction combine these into the desired independence
of \(X_D\) and \(X_E\). The case \(v\in E\) is symmetric.

If \(v\in C\), put \(C'=C\setminus\{v\}\). Parents of a
common child form a clique in the moral graph. Thus its parents
outside \(C'\) cannot meet both \(D\) and \(E\). Suppose they
are all in \(D\); the other case is symmetric. The terminal kernel
gives \(X_v\perp X_E\mid X_D,X_{C'}\), while induction on the
remaining graph gives \(X_D\perp X_E\mid X_{C'}\). Symmetry and
contraction give \((X_D,X_v)\perp X_E\mid X_{C'}\); weak union
then gives \(X_D\perp X_E\mid X_{C'},X_v\). This is the required
conditioning on \(X_C\). The argument also covers the case where
all parents lie in \(C'\). This completes the full-partition induction.

Now let disjoint \(A,B,C\) be d-separated in \(G\). Restrict to
\(H=\operatorname{An}_G(A\cup B\cup C)\). The ancestral marginal
lemma retains its kernels. By the graph-theoretic moral bridge, no
component of the moral graph of \(G[H]\) after deletion of \(C\)
meets both \(A\) and \(B\). Assign all components meeting \(A\)
to \(D\), and all remaining components to \(E\). This fills in
the vertices not originally named in the query and gives a full
partition of \(H\) with no \(D\)-to-\(E\) edge, with
\(A\subseteq D,B\subseteq E\). The induction just proved gives
\(X_D\perp X_E\mid X_C\); decomposition gives the desired query.
Empty query sets are immediate. No density, positivity or intersection
property was used.

\paragraph{Global implies local, and local implies ordered.}
For a vertex \(v\), the set consisting of \(v\) and its
nondescendants is ancestral. Within it \(v\) has no children.
Its only moral neighbors are its parents: any other moral neighbor
would require a common child, which would be a descendant of \(v\).
Deleting the parents therefore separates \(v\) from its other
nondescendants. The global property gives the local property.
Every predecessor is a nondescendant, so decomposition gives the
ordered property.

\paragraph{Ordered implies recursive kernels.}
For each vertex choose a regular conditional kernel of \(X_v\)
given its parents, specified by a measurable probability default on
any null parent arguments. The ordered property and the first
conditional-independence equivalence make this same kernel a
conditional law of \(X_v\) given all predecessors. In detail,
compare its values on a countable generating class with a regular
law given all predecessors; outside their countable null union they
agree as measures. The joint-measure identity for conditional kernels
then expresses each successive actual prefix law as its previous
prefix combined with the chosen kernel. Induction gives the full
recursive kernel law. Together the implications prove all four
equivalences, with the exact standard Borel existence hypotheses.
\end{proof}
\end{annexitem}

\begin{annexitem}{annex:p16.measurable-interventions}{thm:p16-measurable-interventions}{Independent disturbances and measurable equation replacement}
\begin{proof}
Topological induction gives a unique solution at every background
value: each total function is evaluated after its parent values are
known. The same induction, using product measurability of finite
tuples and composition, proves joint measurability in all fixed
intervention values and background coordinates. This is the argument
of \cref{thm:p16-mediation-measurability} with the intervention
parameters retained explicitly.

For an unmodified vertex define
\(K_v(x_{\operatorname{pa}(v)},B)=
\int\mathbf1_B(f_v(x_{\operatorname{pa}(v)},u))\lambda_v(du)\),
where \(\lambda_v\) is its disturbance law. The kernel section
theorem makes this measurable in the parent tuple, and its values
are probability measures. Every predecessor in a topological order
is a function of earlier disturbances, independent of the current
\(U_v\). Successive integration therefore gives exactly this
conditional kernel. Under a constant intervention its output is
the prescribed value and its kernel is \(\delta_{t_v}\). Thus
the full intervened law is the claimed recursive kernel law.
Integrating a measurable indicator of the joint solution against
the fixed product disturbance law makes \(t\mapsto P^t(B)\)
measurable for every target event, proving a probability kernel in
the intervention parameter.

A history-based replacement kernel is realized with a new independent
uniform variable by the randomization theorem. Its parents are the
specified earlier coordinates; with the assumed acyclic augmented
graph, the same argument gives its kernel law and global Markov
property. For consistency, the factual solution already satisfies
all replacement equations when the replaced values are its attained
values, and uniqueness preserves the solution. The same reasoning
in an already modified model proves composition. These are pointwise
equation identities and require no observational support condition.
Changing a kernel on an unobserved parent argument, however, changes
the specified causal model and can change interventions there.
\end{proof}
\end{annexitem}

\begin{annexitem}{annex:p16.measurable-response-component}{thm:p16-measurable-response-component}{A whole response background on one side of the separator}
\begin{proof}
Add the private disturbance roots and take the ancestral graph of
\(Y,B,C\), including those roots, then moralize and delete \(C\).
Let \(D\) be the component containing \(Y\). It contains no vertex
of \(B\), by the separation hypothesis. Every parent of a vertex
in the ancestral graph is also in that graph. If a nonconditioning
parent has a child in \(D\), its direct edge remains in the moral
graph and hence it too lies in \(D\). In particular every private
disturbance needed by an endogenous vertex of \(D\) is among the
roots in \(R\). Parents outside \(D\) can only be in \(C\).

Set these external conditioning inputs to an arbitrary formal tuple
\(c\), and the private roots of \(D\) to an arbitrary \(r\).
Evaluate the vertices of \(D\) in their inherited topological order,
using the given parameter \(t\). Total functions define every
coordinate even if \((c,r)\) is not naturally attainable. Joint
measurability follows by induction. Denote the resulting response
function by \(F(t,c,r)\). At an actual outcome, the conditioning
tuple is unchanged by hypothesis, and all parent and disturbance
inputs agree with those in this restricted evaluation. Induction
therefore proves \(Y^t=F(t,C,R)\) simultaneously, pointwise for
every parameter. No almost-everywhere choice was involved.

Choose one parameter value \(t_0\). In the model at \(t_0\),
the tuple of roots \(R\) is d-separated from \(B\) given \(C\):
their relevant ancestral moral graph is a subgraph of the one already
considered, and its root vertices lie on the \(D\) side. The global
Markov theorem therefore gives \(R\perp B\mid C\). Both these
tuples and their background law are the same for all parameters,
so this is one assertion about their original joint law, not a
parameter-dependent collection of versions. The common product-kernel
characterization supplies one conditional measure identity.

If the random parameter is a measurable function of \((B,C)\), its
conditional joint law with \(R\) is the corresponding pushforward
of this product kernel. Integrating the jointly measurable \(F\)
against that joint law proves the stated random-substitution license,
for bounded tests or with explicit absolute integrability for means.
This is exactly the uniform premise unavailable from separate
pointwise conditional-response assertions.
\end{proof}
\end{annexitem}

\begin{annexitem}{annex:p16.measurable-mediation-fields}{thm:p16-measurable-mediation-fields}{Three uniform response fields from three different graph separations}
\begin{proof}
All the graphs in this proof are finite, and all their functions remain
jointly measurable. Expose independent disturbances as private roots;
this does not change any separation among the original vertices, since
a new root has just one child. Baseline nondescendant conditions ensure
that the tuples named below are unchanged pointwise by the specified
input replacements.

\paragraph{Cross-world field.}
Take \(H=G_{\underline A\underline M}\), restrict to the ancestors
of \(\{Y\}\cup M\cup C\), moralize and remove \(C\). Let \(D_Y\)
be the component of \(Y\) and \(D_M\) the union of components meeting
the mediator set. The first source separation makes them disjoint.
Restore every original outgoing mediator edge whose child is an ancestor
of some mediator in the graph with exposure inputs cut. This restores
the mediator-to-mediator ancestral equations needed for the entire tuple
\(M_b\), while keeping the separation intact. Here are the exact
graph facts used in that claim. From the child of any restored edge,
follow a directed route to its first mediator endpoint. The route already
exists in \(H\), as it uses no outgoing mediator edge before that
endpoint. It avoids \(C\), which has no mediator descendants. Thus
the child lies in \(D_M\). Every other remaining nonbaseline parent
of that child either has an edge already in \(H\), or is itself a
mediator; it too lies in \(D_M\). New moral edges therefore join only
\(D_M\) and possibly deleted baseline vertices. All original ancestors
of the mediator set already belong to the ancestral graph, by stopping
their directed routes at the first mediator. No new ancestral vertex
is required and no new moral connection to \(D_Y\) appears.

Fix \(b\). In the restored graph, replace every exposure input to a
child in \(D_Y\) by \(a\) and every other exposure input by \(b\).
Replace every mediator input not restored by its component of \(m\).
The mediator tuple in this auxiliary model is exactly \(M_b\), by
topological evaluation on its restored ancestors. The conditioning
tuple is the original \(C\). Both are unchanged as \((a,m)\) varies.
On ancestral routes to \(Y\), stop at \(C\). All other vertices
needed there belong to \(D_Y\), and their exposure and mediator inputs
are precisely \(a,m\); hence the response is \(Y_{a,m}\) pointwise.
These equation comparisons use no value-space cardinality assumption.

Apply the response-component lemma to this parameterized auxiliary
model, with parameter \((a,m)\), unchanged tuple \(B=M_b\) and
conditioning \(C\). It produces the actual private-root tuple \(R_0\),
a jointly measurable \(F_0\), and \(R_0\perp M_b\mid C\). The
response component and its equations depend on \(a,m,C,R_0\), not
on \(b\). Thus the same \(R_0,F_0\) works for each reference value
\(b\). This establishes a common product-kernel premise for the whole
response family, rather than only independence of its separately
evaluated fixed-mediator coordinates.

\paragraph{Exposure-adjustment field.}
In \(G_{\underline A\overline M}\), retain the natural exposure
equation, feed \(a\) to its outgoing inputs, and fix all mediator
equations at \(m\). Its response is \(Y_{a,m}\). Its tuples
\(A\) and \(S=(C,W_1)\) are unchanged pointwise: every mediator
is a descendant of \(A\), and \(S\) has no exposure descendants.
The second separation and the response-component lemma, now with
parameter \((a,m)\), unchanged tuple \(A\) and conditioning \(S\),
give \(Y_{a,m}=F_1(a,m,S,R_1)\) for all parameters and
\(R_1\perp A\mid S\). Disturbances at fixed mediator coordinates
can be retained as unused inputs; their extra private roots do not
create a path between old vertices and can be ignored by \(F_1\).

\paragraph{Mediator-adjustment field.}
Start instead from \(G_{\underline M}\) with conditioning tuple
\(T=(A,C,W_1,W_2)\). None of its coordinates is a mediator descendant.
Restore mediator ancestral edges by the preceding restoration argument.
Keep the natural mediator tuple and replace its remaining outgoing
inputs by \(m\). The response is \(Y_m\), while \(M,T\) are
unchanged pointwise as \(m\) varies. The third separation survives
restoration and the response-component lemma gives a jointly measurable
\(F_2\) and actual private-root tuple \(R_2\) satisfying
\(Y_m=F_2(m,T,R_2)\) for every \(m\), with \(R_2\perp M\mid T\).
In particular consistency supplies the simultaneous measurable identity
\(Y=F_2(M,T,R_2)\). This identity, not conditioning on a null equality
event between two continuously distributed mediators, is what will
identify the observed response kernel.

\paragraph{The reference-mediator law.}
Finally retain natural \(A\) and replace its outgoing inputs by
\(b\). The mediator is \(M_b\), and \((C,W_3)\) is unchanged.
The fourth separation and the general global Markov theorem give
\(M_b\perp A\mid C,W_3\). Positive conditional exposure atoms
identify its regular conditional law there with the observed law
of \(M\) in exposure group \(b\), using consistency. The equality
is outside one \((C,W_3)\)-null set as measures: compare a countable
generating class and use measure uniqueness. Integrate over
\(\gamma_3\) to obtain the displayed \(\eta_b\) as the law of
\(M_b\) given \(C\), again with one common baseline-null set.
All three field independences have this same common-measure
interpretation by the conditional-independence calculus.
\end{proof}
\end{annexitem}

\begin{annexitem}{annex:p16.measurable-mediation-corollary}{thm:p16-measurable-mediation-corollary}{Independent reference draw, sequential replacement and overlap transfer}
\begin{proof}
\paragraph{Replace the natural coupling by an independent draw using its proved field.}
Construct a new mediator \(G^b\) with joint law
\(P(d\omega)\eta_b(C(\omega),dg)\), as in the preceding mediation
policy definition. This uses the original background, together with a
fresh draw conditional on baseline. The kernel joint-measure theorem
proves existence. Conditional on \(C\), the joint law of
\((M_b,R_0)\) is the product of \(\eta_b\) and the conditional
\(R_0\) law by the first field. The enlarged model has exactly that
same product law for \((G^b,R_0)\). The single jointly measurable
map \(F_0(a,\cdot,C,\cdot)\) therefore gives
\[
 \mathcal L(Y_{a,M_b})=\mathcal L(Y_{a,G^b}).               \tag{1}
\]
This is a statement about equal pushforwards of equal joint measures;
it involves no pointwise conditional-version substitution.

\paragraph{Use the exposure field on that independent draw.}
Because \(S\) contains \(C\), the conditional law of
\((R_1,A,G^b)\) given \(S\) factors as
\[
 \mathcal L(R_1\mid S)\otimes\mathcal L(A\mid S)
                         \otimes\eta_b(C,\cdot).
\]
To verify this, integrate a rectangle indicator under the enlarged
measure. The new draw contributes the factor \(\eta_b(C,D)\),
which is measurable from \(S\), and the first two factors follow
from \(R_1\perp A\mid S\). A countable generating-class argument
gives equality as kernels outside one \(S\)-null set. Consequently
\((R_1,G^b)\perp A\mid S\), and the joint response family gives
\(Y_{a,G^b}=F_1(a,G^b,S,R_1)\perp A\mid S\).
For a bounded \(\varphi\), positive \(P(A=a\mid S)\) now permits
\[
 E[\varphi(Y_{a,G^b})]
 =\int E[\varphi(Y_{a,G^b})\mid A=a,S=s]\,P_S(ds).         \tag{2}
\]
The baseline law of \(S=(C,W_1)\) is \(\nu\gamma_1\).
Its law given \(A=a\) has density
\(P(A=a\mid S)/P(A=a)\) with respect to it; this is a positive
Radon--Nikodym derivative almost everywhere, not an assumed coordinate
density. Thus the two laws are equivalent, and the conditional equality
in (2) holds on every needed measure class.

\paragraph{Use the mediator field before evaluating its random argument.}
Let \(\kappa_2(t,dr)\) be a conditional law of \(R_2\) given
\(T=t\). The third field gives one product kernel for
\((M,R_2)\) given \(T\). Since
\(Y=F_2(M,T,R_2)\) pointwise, the jointly measurable kernel
\[
 \widetilde Q(t,m,B)=\int\mathbf1_B(F_2(m,t,r))\kappa_2(t,dr)
\]
is a conditional law of the observed \(Y\) given \((T,M)\).
This follows by integrating the actual product law on rectangles;
kernel uniqueness extends the equality to all response events on
one \(P_{T,M}\)-null set. In exposure group \(a\), it therefore
agrees with the specified \(Q_a^{\mathrm{obs}}\) for the observed
law of \((C,W_1,W_2,M)\) in that group.

The new draw has law \(\eta_b(C,\cdot)\) independent of \(R_2\)
conditional on \(T\): the same rectangle calculation applies because
\(T\) includes \(C\). Thus a conditional law of \(Y_{G^b}\) given
\(T=t\) is \(\int\widetilde Q(t,m,\cdot)\eta_b(c,dm)\).
On the event \(A=a\), intervention on \(M\) leaves \(A\) unchanged,
since every mediator is its descendant. Composition at that attained
exposure value gives \(Y_{a,G^b}=Y_{G^b}\) pointwise in the enlarged
space. Partitioning the right side of (2) by \(W_2\) consequently
gives the target integral with \(\widetilde Q\), \(\eta_b\),
\(\gamma_2^a\), \(\gamma_1\), and \(\nu\).

It remains to justify replacement of \(\widetilde Q\) by the
\emph{observed} response kernel on this new integration measure.
Let \(\alpha=\nu\gamma_1\gamma_2^a\). The observed group law
of \((C,W_1,W_2,M)\) is equivalent to \(\alpha\rho_a\),
by the positive exposure derivative above. If a measurable set has
zero measure under this law, nonnegative kernel integration implies
that its mediator section has \(\rho_a\)-mass zero for
\(\alpha\)-almost every \((c,w_1,w_2)\). Condition (MS) gives
zero \(\eta_b\)-mass for all those sections. Therefore the
target measure \(\alpha\eta_b\) is absolutely continuous with
respect to the observed group law. The single exceptional set on
which the two conditional response measures differ has target mass
zero. Their integrals are identical, proving the stated observable
formula together with (1).

For means, first apply the bounded formula to \(|y|\wedge n\).
Monotone convergence identifies the target absolute integral with
the absolute expectation of the nested response. Under its stated
finiteness, positive and negative parts give the signed formula.
The observed integrability condition permits using the ordinary
conditional mean version, with a measurable default on the
observed-null infinite-integral set; absolute continuity also removes
that set from the target. Indicator tests identify the entire response
law. Applying this result to a finite collection of exposure pairs
and taking a finite union of the relevant null sets licenses all
the stated effect subtractions.
\end{proof}
\end{annexitem}

\begin{annexitem}{annex:p16.measurable-graph-examples}{thm:p16-measurable-graph-examples}{Singularity does not prevent kernel factorization}
\begin{proof}
The chain model has kernels \(K_X\) uniform, \(K_Z(x,\cdot)=\delta_x\),
and \(K_Y(z,\cdot)=\delta_z\). They are measurable probability
kernels, so the global Markov theorem applies. Its law gives
probability one to \(D=\{(x,z,y):x=z=y\}\). Suppose a product
\(\mu_1\otimes\mu_2\otimes\mu_3\) of sigma-finite measures
dominated this law. A sigma-finite measure on a real Borel space has
at most countably many positive-mass singletons: cover the space by
countably many finite-measure sets, and in each such set there are
only finitely many singletons with mass at least \(1/n\), for each
positive integer \(n\). Let \(J\) be the countable set of points with
positive \(\mu_2\)-mass. Every singleton has finite mass under a sigma-finite
measure, since it lies in one member of such a finite-measure cover.
For \(x\notin J\), the section of
\(D\cap\{x\notin J\}\) has product mass
\(\mu_2(\{x\})\mu_3(\{x\})=0\). Tonelli therefore gives this
whole diagonal subset product measure zero. Its actual probability
is one, since a uniform variable avoids the countable \(J\).
This contradicts absolute continuity. Thus a product-density proof
cannot be presumed for these kernels.

In the second model, \(P(M_0\leq m)=m\) and
\(P(M_1\leq m)=P(U\leq m^2)=m^2\), for \(0<m<1\).
The corresponding densities are \(1\) and \(2m\), and both
are positive on the mediator interval. They are mutually absolutely
continuous, so (MS) holds in both directions. Both exposure atoms
are \(1/2\). In the DAG \(A\to M\to Y\), \(A\to Y\),
all four required separations hold with empty covariate sets:
cutting outgoing exposure edges isolates exposure; cutting outgoing
mediator edges leaves only the fork through conditioned exposure;
cutting both removes all connections between mediator and response.
Independent mediator noise and a constant response disturbance give
the required structural model.

The controlled responses are \(Y_{0,m}=m\) and \(Y_{1,m}=1+2m\).
Their integrals against \(1\,dm\) and \(2m\,dm\) on \((0,1)\)
give respectively
\[
 \int_0^1m\,dm=1/2,\qquad \int_0^12m^2\,dm=2/3,
 \qquad1+2(1/2)=2,\qquad1+2(2/3)=7/3.
\]
Subtraction gives PNDE \(3/2\), TNIE \(1/3\), and TE \(11/6\).
The response is on the curve \(y=a+m+am\) rather than having
an ordinary conditional density on a response interval, but the
conditional response probability kernel is its point mass and
the proof remains valid.

For completeness, positive exposure density does not by itself
identify a fixed null exposure point. Let \(A,M\) be independent
uniform variables on \((0,1)\), and fix \(a_*=1/2,b_*=1/4\).
Take also an independent Bernoulli\((1/4)\) bit \(E\), and compare
\(Y=E\oplus\mathbf1_{\{A=a_*\}}\) with \(Y=E\).
Permit the same DAG \(A\to M\to Y,A\to Y\), with
unused permitted inputs. Both models have independent disturbances,
identical observed law \((A,M,E)\), the four empty-set graph
separations, and identical positive mediator densities. Yet their
natural direct effects for \(b_*\) to \(a_*\) are \(3/4-1/4=1/2\)
and zero. The observed density with respect to Lebesgue measure in
\((A,M)\) and counting measure on the binary response is \(3/4\)
at response zero and \(1/4\) at response one, everywhere on the
open unit square; it is strictly positive.
The structural intervention is defined in either model, but a
conditional response value at the null exposure is not identified
by the observed law. Thus the conditional-atom domain in the
observational corollary is a stated sufficient domain; replacing
it by an unqualified point-density condition would be false.
\end{proof}
\end{annexitem}
